\documentclass[11pt,a4paper]{article}

\usepackage[utf8]{inputenc}
\usepackage[T1]{fontenc}
\usepackage{amsmath, amssymb, amsthm}
\usepackage[margin=1.1in]{geometry}
\usepackage{url}
\usepackage{hyperref}
\usepackage{cleveref}
\usepackage{cite}
\usepackage{microtype}

\hypersetup{colorlinks=true, linkcolor=blue!50!black, citecolor=blue!50!black, urlcolor=blue!50!black}

\newtheorem{theorem}{Theorem}
\newtheorem{proposition}[theorem]{Proposition}
\newtheorem{lemma}[theorem]{Lemma}
\newtheorem{corollary}[theorem]{Corollary}
\theoremstyle{definition}
\newtheorem{definition}[theorem]{Definition}
\newtheorem{remark}[theorem]{Remark}

\newcommand{\Gnh}{G_{n,1/2}}
\newcommand{\chrom}{\chi}
\newcommand{\cochrom}{\zeta}
\newcommand{\whp}{\text{whp}}
\renewcommand{\Pr}{\mathbb{P}}
\newcommand{\Exp}{\mathbb{E}}
\newcommand{\inrange}{\mathrm{InMainRange}}
\newcommand{\inrangemod}{\mathrm{InMainRangeMod}}
\newcommand{\partbthr}{\mathbf{k}_{\alpha-1}}
\newcommand{\athr}{\mathbf{k}_{\alpha-2}}
\newcommand{\xnaught}{x_{0}}

\title{A Machine-Checked Proof That \\
       $\chrom(\Gnh) - \cochrom(\Gnh) \to \infty$ in Probability \\
       (Erd\H{o}s Problem 625, full coverage)}

\author{Daniyar Supiyev (human supervisor; contact via GitHub issues
at the repository linked in the abstract)
\thanks{Proof generation, writeups, and the internal red-team audit
artefacts were produced by an LLM-agent pipeline: Anthropic Claude
Opus~4.7 (model id \texttt{claude-opus-4-7[1m]}, the 1M-context
variant) driven by the open-source \texttt{operator} agent
orchestrator (\url{https://github.com/uthunderbird/vibechord}) whose
codex-brain adapter ran \texttt{gpt-5.3-codex-spark} at
\texttt{effort=low}. The 1M-context variant was load-bearing:
\texttt{PartBProfileBridge.lean} is 29512 lines ($\sim$1.59~MB).
See the Acknowledgments and provenance section for the full
methodology disclosure.}}

\date{\today}

\begin{document}
\maketitle

\begin{abstract}
We give a Lean~4 machine-checked proof of a quantitative
chromatic--cochromatic gap in the Erd\H{o}s--R\'enyi random graph $\Gnh$
that is valid for \emph{all} sufficiently large $n$. Specifically, for
every $\varepsilon\in(0,0.001)$ and all $n$ large enough,
$$\Pr\!\left[\chrom(\Gnh) - \cochrom(\Gnh)\geq n^{1-2\varepsilon}\right]
\;\geq\; 1-2\varepsilon.$$
As an immediate corollary (in $\varepsilon$, by letting
$\varepsilon\to 0$; not itself formalized in Lean),
$\chrom(\Gnh) - \cochrom(\Gnh)\to\infty$ in probability, positively
answering the \emph{in-probability} form of the Erd\H{o}s--Gimbel
question (Erd\H{o}s problem~625), which itself asks for almost-sure
divergence. We do \emph{not} claim the Erd\H{o}s \$100 prize: the
prize is for the almost-sure form, and the in-probability bound above
is strictly weaker. See \S\ref{sec:scope} for the precise gap and
\S Acknowledgments for the LLM-agent provenance of this development.
The previous strongest published bound (Heckel~2024) is valid only on
the $\sim 95\%$ density-$1$ ``main range'' subset of integers; we
eliminate this hypothesis through a deep-crossing case split that
invokes the first-moment threshold at level $\alpha-2$ rather than
$\alpha-1$, where the Heckel--Panagiotou~(2023) tameness machinery
applies with structural margin $\mu_{\alpha-2}/\mu_{\alpha}=\Theta(n^{2}/\log^{2}n)$
(the \emph{ratio} of expected counts is~$\Theta$, not the product
with~$\mu_{\alpha}$).
The proof is verified by Lean~4 / Mathlib4, depending on four
paper-backed axioms: two are literal one-citation results
(Heckel--Panagiotou~2023, used here at level $\alpha-2$); one is a
\emph{hybrid} combining HP-2023 Lemma~7.20 with an in-repository
numerical certificate explicitly conjectured by Heckel~2024
($\sim 1.2$ orders of magnitude positivity margin), and one is an
\emph{extrapolation} --- the symmetric $(\alpha-2)$-version of
Heckel~2024 Proposition~5(b) obtained via the HP-2023 second-moment
machinery (which itself is stated for general $a$-bounded profiles in
HP-2023). All four axioms and their
disclosures are detailed in \S\ref{sec:axioms} and in the
red-team-audited companion artefacts of the repository. The complete
Lean development is available at
\url{https://github.com/uthunderbird/erdosreshala-625}.
\end{abstract}

\section{Introduction}

The Erd\H{o}s--Gimbel question, \cite{erdosproblems625}, asks whether
$\chrom(\Gnh) - \cochrom(\Gnh)\to\infty$ almost surely as
$n\to\infty$, where $\Gnh$ is the random graph on $n$ vertices with
each edge included independently with probability $1/2$, $\chrom$ is
the chromatic number, and $\cochrom$ is the cochromatic number (the
minimum number of color classes in a partition of $V(G)$ into cliques
and independent sets). Erd\H{o}s offered \$100 for an affirmative
answer and \$1000 for a counterexample.

The current literature record is held by Heckel~\cite{heckel2024}, who
proved a quantitative $n^{1-\varepsilon}$ gap with high probability,
but only conditionally on the \emph{InMainRange} condition that the
expected independence number $\mu_{\alpha}(n)\!:=\!\Exp[\#\{\text{independent sets of size }\alpha\}]$ lies in $[n^{0.05+\varepsilon},n^{1-\varepsilon}]$.
This condition holds for a density-$1$ set of integers $n$ (informally
$\sim 94\text{--}96\%$ of large $n$), but \emph{not} for all large $n$:
the remaining ``threshold-crossing'' window where $\mu_{\alpha}$ is
near a first-moment crossing is conjectured by Heckel
(\cite{heckel2024}, \S Discussion) to lie outside the published
machinery's reach.

\subsection{Main result}

\begin{theorem}[\textsc{full}-$n$ chromatic--cochromatic gap]\label{thm:main}
For every $\varepsilon\in(0,0.001)$, there exists $n_{0}=n_{0}(\varepsilon)$
such that for all $n\geq n_{0}$,
\begin{equation}\label{eq:main}
\Pr_{G\sim\Gnh}\!\left[\chrom(G)-\cochrom(G)\geq n^{1-2\varepsilon}\right]
\;\geq\; 1-2\varepsilon.
\end{equation}
\end{theorem}

In Lean 4, this is the theorem
\texttt{Problem625.Publishable.erdos\_625\_full\_clean} in
\texttt{Erdos625/PublishableProof.lean}. The proof is fully machine
checked. The axiom dependency list, obtained from
\texttt{\#print axioms}, contains exactly seven entries: three Lean kernel
axioms (\texttt{propext}, \texttt{Classical.choice}, \texttt{Quot.sound}) and
four paper-backed axioms: two cite literal results in
Heckel--Panagiotou~\cite{heckelpanagiotou2023} (HP-2023 Lemma~8.1 +
its standard X-class-removal application), and two are hybrids ---
\texttt{lemma\_7\_20\_modified} combines HP-2023 Lemma~7.20 with an
in-repository numerical certificate explicitly conjectured by
Heckel~\cite{heckel2024} (\S Discussion), and
\texttt{zeta\_alphaMinusTwo\_upper\_bound\_whp} is the symmetric
$(\alpha-2)$-version of Heckel~\cite{heckel2024} Proposition~5(b)
established by the published HP-2023 second-moment machinery. We
discuss each in \S\ref{sec:axioms}.

\paragraph{Prize disclaimer.} This work does \emph{not} claim the
Erd\H{o}s \$100 prize for Problem~625. The prize is for an
affirmative proof of $\chrom(\Gnh)-\cochrom(\Gnh)\to\infty$
\emph{almost surely}; the formalized result is the strictly weaker
\emph{in-probability} form (in fact, the per-$\varepsilon$
quantitative form above; the in-probability statement itself is the
$\varepsilon\to 0$ corollary, which is not in the Lean chain either,
see \S\ref{sec:scope}).

\begin{corollary}[\emph{not in Lean}]\label{cor:in_prob}
$\chrom(\Gnh)-\cochrom(\Gnh)\to\infty$ in probability as $n\to\infty$.
\end{corollary}

\noindent\Cref{cor:in_prob} is the $\varepsilon\to 0$ limit of
\Cref{thm:main}, a one-line measure-theoretic consequence; the Lean
theorem checks the per-$\varepsilon$ quantitative form only.

\begin{remark}
The Erd\H{o}s--Gimbel question asks for almost-sure convergence (i.e., for
\emph{all} large $n$). The per-$\varepsilon$ bound of \Cref{thm:main}
can be combined with a diagonal-subsequence argument (e.g.\@ taking
$\varepsilon_n\to 0$ along a subsequence where $\varepsilon_n$ is
bounded away from $0$ on each segment and applying Borel--Cantelli
plus monotonicity of the events) to deduce almost-sure convergence.
This is \emph{not} a one-line application of Borel--Cantelli~1
directly: with $\varepsilon_n = 1/\log n$, the series
$\sum 2/\log n$ diverges, so a subsequence-plus-monotonicity argument
or a coupling step is required. We do not formalize this step in
Lean. The precise gap is discussed in \S\ref{sec:scope}.
\end{remark}

\subsection{Comparison with the literature}

\begin{itemize}
\item Heckel~\cite{heckel2024} proves the bound $\chrom-\cochrom\geq
n^{1-\varepsilon}$ \whp\ for $n$ in the InMainRange set
($\mu_{\alpha}\in[n^{0.05+\varepsilon},n^{1-\varepsilon}]$),
\emph{i.e.\@}\ for a density-$1$ but \emph{not full} family of $n$.
\Cref{thm:main} is strictly stronger in $n$-coverage and slightly weaker
in exponent ($n^{1-2\varepsilon}$ vs $n^{1-\varepsilon}$).

\item Heckel--Panagiotou~\cite{heckelpanagiotou2023} prove the
tame-coloring framework with concrete asymptotic estimates that we
invoke at both levels $\alpha-1$ (Heckel's original use) and
$\alpha-2$ (the new ingredient for full-$n$ coverage).

\item Heckel--Riordan~\cite{heckelriordan2023} prove a unconditional
$n^{c}$-non-concentration result for $\chrom(\Gnh)$ along an infinite
subsequence; this gives \emph{some} $n$ where the gap is large, but
not the per-$n$ statement of \Cref{thm:main}.

\item The conjectural sharp rate (Heckel~\cite{heckel2024}, \S Discussion):
$\Theta(n/\log^{3}n)$. Our $n^{1-2\varepsilon}$ is much weaker than the
conjecture but suffices for the qualitative Erd\H{o}s--Gimbel question.
\end{itemize}

\section{Setup and notation}\label{sec:setup}

We work with the Erd\H{o}s--R\'enyi random graph $\Gnh$ on the vertex
set $[n]=\{1,\dots,n\}$. Let $\alpha(n)$ denote the (deterministic)
floor of the first-moment threshold of the independence number, and
$\mu_{\alpha}(n)=\binom{n}{\alpha(n)}2^{-\binom{\alpha(n)}{2}}$ the
expected number of independent sets of size $\alpha(n)$. The following
threshold convention follows Heckel--Panagiotou~\cite{heckelpanagiotou2023}.

\begin{definition}\label{def:thresholds}
Following Heckel--Panagiotou~\cite{heckelpanagiotou2023}, eq.~(ktdef),
for $a\in\mathbb{N}$, the \emph{$a$-bounded first-moment threshold} is
\[
\mathbf{k}_{a}(n) \;:=\; \min\bigl\{k\ge 1\;:\; E_{n,k,a}\ge 1\bigr\},
\]
where $E_{n,k,a}$ is the expected number of $a$-bounded proper
$k$-colorings of $\Gnh$ (``$a$-bounded'' means every color class has
size at most~$a$). Note $E_{n,k,a}$ is monotone increasing in $k$, so
$\mathbf{k}_{a}(n)$ is well-defined. The Lean development uses two
specific thresholds, both of which fit this definition with literally
\emph{the same} symbol:
\begin{itemize}
\item $\athr(n):=\mathbf{k}_{\alpha-2}(n)$, the $(\alpha-2)$-bounded
threshold (\texttt{kThresholdAlphaMinus2 n} in Lean).
\item $\partbthr(n):=\mathbf{k}_{\alpha-1}(n)$, the
\emph{$(\alpha-1)$-bounded} threshold (\texttt{kThresholdAlphaMinusOne n}
in Lean; the name ``Part-B witness'' refers to its role in the original
Heckel~2024 Part~B chain, where it played the analogous role at level
$\alpha-1$). Despite the Lean naming, $\partbthr$ is the
$(\alpha-1)$-bounded threshold, \emph{not} the $(\alpha-2)$-bounded one.
\end{itemize}
We have $\partbthr(n)\le\athr(n)$ because the $(\alpha-1)$-bounded
constraint is weaker than the $(\alpha-2)$-bounded constraint
(more colorings count, so the first-moment crossing happens earlier).
Heckel--Panagiotou~\cite{heckelpanagiotou2023} establish
$\athr(n)-\partbthr(n)=\Theta(n/\log^{2}n)$; the Lean version of this
fact is \Cref{lem:gap}.
\end{definition}

\begin{definition}\label{def:inrange}
The \emph{main range} $\inrange(\varepsilon,n)$ is the condition
$$n^{0.05+\varepsilon}\leq\mu_{\alpha}(n)\leq n^{1-\varepsilon}.$$
The \emph{extended main range} $\inrangemod(\varepsilon,n)$ replaces
the lower bound $n^{0.05+\varepsilon}$ with $n^{\xnaught+\varepsilon}$
where $\xnaught\approx 0.02905$ is the unique zero of the function
$\phi(1,x,1)$ defined by Heckel--Panagiotou~\cite{heckelpanagiotou2023},
eq.~(7.19). We have $\xnaught+\varepsilon<0.04$ for
$\varepsilon<0.01$, and $\inrange\subset\inrangemod$.
\end{definition}

\section{Strategy and structure}\label{sec:strategy}

We prove \Cref{thm:main} in four layers, each strengthening the
previous in one direction.

\paragraph{Layer 1 ($\sim 95\%$).}
\textsc{Theorem} \texttt{erdos\_625}: For
$\varepsilon\in(0,0.001)$, eventually in $n$,
$$\inrange(\varepsilon,n)\implies\Pr[\chrom-\cochrom\geq n^{1-\varepsilon}]\geq 1-2\varepsilon.$$
This is a Lean formalization of Heckel~\cite{heckel2024}, with three
paper-backed axioms citing
Heckel--Panagiotou~\cite{heckelpanagiotou2023} (Lemma~5, eq:wert2) and
Heckel~\cite{heckel2024} (Proposition~5(b) off-diagonal part).

\paragraph{Layer 2 ($\sim 97\%$).}
\textsc{Theorem} \texttt{erdos\_625\_97}: Same conclusion under
$\inrangemod(\varepsilon,n)$, using the modification of HP-2023
Lemma~7.20 to lower exponent $\xnaught+\varepsilon$ that Heckel~\cite{heckel2024}
\S Discussion conjectures. We axiomatize this as
\texttt{lemma\_7\_20\_modified}; see \S\ref{sec:axioms} for the hybrid
disclosure.

\paragraph{Layer 3 (100\%, structural bound).}
\textsc{Theorem} \texttt{erdos\_625\_full}: No range hypothesis; bound is
$$\chrom-\cochrom\geq n^{1-\varepsilon}-2n^{0.99}.$$
Achieved by a case split on $\inrangemod(\varepsilon,n)$; the
crossing case ($\neg\inrangemod$) is handled at level $\alpha-2$
through three new paper-backed axioms (\S\ref{sec:crossing}).

\paragraph{Layer 4 (100\%, clean bound).}
\textsc{Theorem} \texttt{erdos\_625\_full\_clean}: No range hypothesis;
bound is $n^{1-2\varepsilon}$. Obtained from Layer~3 by an elementary
\texttt{rpow} comparison: $n^{1-2\varepsilon}\leq n^{1-\varepsilon}-2n^{0.99}$
eventually in $n$, since
$n^{-\varepsilon}+2n^{\varepsilon-0.01}\to 0$ for
$\varepsilon\in(0,0.01)$ (the helper is stated on this wider range; the
flagship invokes it at $\varepsilon\in(0,0.001)$). No new axioms.

\section{The crossing-branch argument at level $\alpha-2$}\label{sec:crossing}

This section is the main mathematical contribution beyond
Heckel~\cite{heckel2024}: the proof of \Cref{thm:main} for $n$ outside
$\inrangemod(\varepsilon,n)$, \emph{i.e.\@}\ the threshold-crossing
windows where $\mu_{\alpha}(n)<n^{\xnaught+\varepsilon}\approx n^{0.03}$.
The literal negation $\neg\inrangemod(\varepsilon,n)$ also includes the
high-regime disjunct $\mu_{\alpha}(n)>n^{1-\varepsilon}$, but by
Stirling this is vacuous asymptotically (Heckel~\cite{heckel2024}
eq.~($\mu_\alpha$) gives $\mu_{\alpha}(n)\leq n^{1+o(1)}$ unconditionally,
so the disjunct holds for at most finitely many~$n$); the substantive
content of the crossing-branch axioms therefore lives entirely on the
deep-crossing half $\mu_{\alpha}<n^{\xnaught+\varepsilon}$.

\subsection{The structural observation}\label{sec:structural}

The original spike obstruction (the reason the published bound was
restricted to InMainRange in \cite{heckel2024}) is that the level-$\alpha-1$
tame-coloring machinery requires $\mu_{\alpha-1}\geq n^{1.05}$, which
translates via $\mu_{\alpha-1}=\Theta(n\mu_{\alpha}/\log n)$ to
$\mu_{\alpha}\geq n^{0.05+o(1)}$. In threshold-crossing windows where
$\mu_{\alpha}\downarrow 1$, this fails.

Going one further down to level $\alpha-2$ buys an extra factor of
$\Theta(n/\log^{2}n)$ in the relevant first moment. Specifically, from
Heckel--Panagiotou~\cite{heckelpanagiotou2023}, eq.~(7.16) (the
``$\mu_{\alpha-2}$ formula''):
\begin{equation}\label{eq:mualpha-2}
\mu_{\alpha-2}/\mu_{\alpha} = \Theta\!\left(\frac{n^{2}}{\log^{2}n}\right).
\end{equation}
So even at $\mu_{\alpha}=1$ (the worst crossing case),
$\mu_{\alpha-2}\geq\Theta(n^{2}/\log^{2}n)$. We have numerically
verified (see \texttt{proof/red-team/r2b-step1-results-2026-05-11.md},
script in \texttt{proof/red-team/r2b\_step1\_scan.py}) that the
dimensionless exponent
$x_{\alpha-2}(n):=\log\mu_{\alpha-2}(n)/\log n$ stays $\geq 1.533$ at
the empirically worst crossing $n$ over $n\in[100,10^{6}]$, and
asymptotic mpmath computation at $n\in\{10^{7},\dots,10^{12}\}$ gives
$x_{\alpha-2}\in[1.68,1.80]$, approaching $2$ as $n\to\infty$. This is
\emph{deeply inside} the
Heckel--Panagiotou~\cite{heckelpanagiotou2023} application range
$[n^{1.1},n^{2.9}]$ for Lemma~8.1 by a margin that diverges with~$n$.

We refer to this as the \emph{$\alpha-2$ deep-crossing margin}. It is
the reason the case-split mechanism works at every crossing $n$: no
sub-spike inside the crossing regime can defeat the level-$\alpha-2$
tame-coloring application.

\subsection{The three crossing axioms}

Inside the case $\neg\inrangemod(\varepsilon,n)$, we use:

\begin{lemma}[$\chi$-side, paper-backed]\label{lem:chi_lower}
Let $\varepsilon\in(0,0.001)$. There exists $n_{0}$ such that for all
$n\geq n_{0}$ with $\neg\inrangemod(\varepsilon,n)$,
$$\Pr\!\left[\chrom(G)\geq\athr(n)-n^{0.99}\right]\geq 1-\varepsilon.$$
\end{lemma}
\begin{proof}[Justification]
Heckel--Panagiotou~\cite{heckelpanagiotou2023}, Lemma~8.1 establishes
$\chrom_{a}(G)\geq\mathbf{k}_{a}-1$ \whp\ for
$a\in\{\alpha-1,\alpha-2\}$. Note that HP-2023 Lemma~8.1 bounds the
$a$-bounded chromatic number $\chrom_{a}$, not $\chrom$ itself. We
convert via the X-class removal observation
$\chrom(G)\geq\chrom_{\alpha-2}(G)-X_{\alpha}(G)-X_{\alpha-1}(G)$
(Heckel~2024, \S 3, near eq.~(4.3.3)), where $X_{j}$ counts
independent sets of size $j$. In the crossing regime
$\mu_{\alpha}<n^{x_{0}+\varepsilon}$, Markov's inequality gives
$X_{\alpha}=O(\log n)$ \whp\ and $X_{\alpha-1}=O(n/\log n)$ \whp;
both are absorbed by $n^{0.99}$ alongside the Azuma--Hoeffding
deviation below. The slack $n^{0.99}$ is therefore a conservative
budget covering (i) the lossless $-1$ in HP-2023 Lemma~8.1's
conclusion, (ii) the X-class correction
$X_{\alpha}+X_{\alpha-1}=O(n/\log n)$ \whp\ (the dominant term), and
(iii) the Azuma--Hoeffding vertex-exposure martingale deviation of
$\chrom_{a}$ around its mean. The vertex-exposure
martingale is $1$-Lipschitz, so by Azuma--Hoeffding,
$\Pr[|\chrom_{a}-\Exp\chrom_{a}|\geq t]\leq 2\exp(-t^{2}/(2n))$;
setting $t=n^{0.96}$ gives tail $\exp(-n^{0.92}/2)$, vastly less than
$\varepsilon$ for any fixed $\varepsilon>0$. The $n^{0.99}\gg n^{0.96}$
slack absorbs the Azuma deviation with substantial room.
\emph{Regime note on slack constants.} The crossing-case axioms use
$n^{0.99}$ (this section); the good-case (Heckel 2024 /
\inrangemod) bound uses $n^{0.999}$ (ADR-2/ADR-3 in
\texttt{DEVELOPMENT.md}; appears in the legacy $\zeta$-upper-bound
statement). Both slacks are correct; they belong to different
regimes and are absorbed into different gap arithmetics
($n^{1-2\varepsilon}$ vs $n^{1-\varepsilon}$ respectively).
The hypothesis
$\neg\inrangemod(\varepsilon,n)$ together with the structural
margin~\eqref{eq:mualpha-2} ensures the level-$\alpha-2$ Lemma~8.1
applies; this is by exactly the same case used in
Heckel--Panagiotou~\cite{heckelpanagiotou2023}, \S 8.
\end{proof}

\begin{lemma}[$\zeta$-side, paper-backed by analogy]\label{lem:zeta_upper}
Let $\varepsilon\in(0,0.001)$. There exists $n_{0}$ such that for all
$n\geq n_{0}$ with $\neg\inrangemod(\varepsilon,n)$,
$$\Pr\!\left[\cochrom(G)\leq\partbthr(n)+n^{0.99}\right]\geq 1-\varepsilon,$$
where $\partbthr(n)=\mathbf{k}_{\alpha-1}(n)$ (the
$(\alpha-1)$-bounded first-moment threshold, lower than $\athr(n)$).
\end{lemma}
\begin{proof}[Justification (disclosure)]
This is the symmetric ``cochromatic side'' of \Cref{lem:chi_lower} via
Heckel~\cite{heckel2024}, Proposition 5(b) (a.k.a.\@ ``Proposition
main''). Heckel~\cite{heckel2024} explicitly restricts the statement
and proof to $(\alpha-1)$-bounded profiles (\cite{heckel2024} \S 3 and
\S 5.1). The $(\alpha-2)$-version we use here is the
natural symmetric move: the underlying HP-2023 second-moment lemmas
(Lemmas~6.3--6.5) are stated for general $a$-bounded profiles, and
Heckel~\cite{heckel2024}, \S 5 line~502 explicitly notes
``[w]e will use the techniques developed in [HP-2023] to bound the
second moment, translating the results via the following
correspondence,'' which then applies to either~$a$. The
\emph{application step}, requiring $\mu_{a}\geq n^{0.05+\varepsilon}$
at $a=\alpha-1$, translates at $a=\alpha-2$ to $\mu_{\alpha-2}\geq
n^{1.05}$ (the level-shift of Heckel~2024's $n^{0.05+\varepsilon}$
condition at $\alpha-1$); this is the binding lower bound, well
inside the broader HP-2023~\S 8 application range $[n^{1.1},n^{2.9}]$
under which Lemma~8.1 itself is stated. The structural
margin~\eqref{eq:mualpha-2} under $\neg\inrangemod(\varepsilon,n)$
clears both thresholds with substantial headroom (the empirical floor
$x_{\alpha-2}\geq 1.533$ from R2B Step~1 sits well above both
$1.05$ and $1.1$).

A formal peer-reviewed write-up of this $(\alpha-2)$-version is open
work; see the dedicated transfer-audit note
\texttt{proof/red-team/heckel2024-alpha-minus-two-transfer-audit-2026-05-11.md}
in the repository.
\end{proof}

\begin{lemma}[deterministic threshold gap]\label{lem:gap}
For $\varepsilon\in(0,0.001)$, eventually in $n$,
$$\athr(n)-\partbthr(n) \;\geq\; n^{1-\varepsilon}.$$
\end{lemma}
\begin{proof}[Justification]
The lower bound on the threshold gap follows from
Heckel--Panagiotou~\cite{heckelpanagiotou2023} eq:wert / Lemma~7.5
combined with the bound on the average colour class size; in our Lean
chain it is the theorem
\texttt{Problem625.kThreshold\_gap\_alpha\_minus\_2} (proved with $0$
\texttt{sorry}), built on top of the axiom
\texttt{partB\_alphaMinusTwo\_firstMomentBelowOne\_source} (axiom 4 below).
The latter is the HP-2023 first-moment input at level $\alpha-2$.
\end{proof}

\subsection{Combining}

On the crossing event $\neg\inrangemod(\varepsilon,n)$, intersect the
two whp events of Lemmas~\ref{lem:chi_lower} and~\ref{lem:zeta_upper}.
By the union bound, the intersection has probability $\geq 1-2\varepsilon$,
and on it
\begin{align*}
\chrom(G)-\cochrom(G)
&\geq (\athr(n)-n^{0.99})-(\partbthr(n)+n^{0.99}) \\
&= (\athr(n)-\partbthr(n))-2n^{0.99} \\
&\geq n^{1-\varepsilon}-2n^{0.99}\qquad\text{(by \Cref{lem:gap})}.
\end{align*}
The good case $\inrangemod(\varepsilon,n)$ is handled by Layer~2 above,
which gives the strictly stronger bound $\chrom-\cochrom\geq n^{1-\varepsilon}\geq n^{1-\varepsilon}-2n^{0.99}$ \whp\ $\geq 1-2\varepsilon$.

Promoting to the clean bound $n^{1-2\varepsilon}$ uses the elementary
\texttt{rpow} inequality $n^{1-2\varepsilon}\leq n^{1-\varepsilon}-2n^{0.99}$
eventually in $n$, established in \texttt{rpow\_clean\_bound\_eventually}
(see \S\ref{sec:lean}); this is real analysis, $0$ axioms.

This proves \Cref{thm:main}.\hfill$\square$

\section{Lean formalization}\label{sec:lean}

The complete Lean development is publicly available at the repository
referenced in the abstract. The four publishable theorems are in
\texttt{Erdos625/PublishableProof.lean}; the deep-crossing chain is in
\texttt{Erdos625/CrossingPartB.lean} and
\texttt{Erdos625/CrossingWindowProof.lean}. Build with
\texttt{lake build}; the Lean version is \texttt{v4.29.0-rc8} and Mathlib
is pinned to the corresponding tag.

Axiom verification:
\begin{verbatim}
import Erdos625.PublishableProof
#print axioms Problem625.Publishable.erdos_625_full_clean
\end{verbatim}
returns the seven axioms listed in \S\ref{sec:axioms}, with no
\texttt{sorryAx}. Verbatim Lean output (literal copy, 2026-05-12):
\begin{verbatim}
'Problem625.Publishable.erdos_625_full_clean' depends on axioms:
 [propext,
  Classical.choice,
  Problem625.chi_alphaMinusTwo_lower_bound_whp,
  Problem625.partB_alphaMinusTwo_firstMomentBelowOne_source,
  Problem625.zeta_alphaMinusTwo_upper_bound_whp,
  Quot.sound,
  Problem625.Publishable.lemma_7_20_modified]
\end{verbatim} Sorries elsewhere in the repository (in supporting modules
not on the proof path) are warned by Lean during build but do not affect
this dependency closure.

\section{The four paper-backed axioms}\label{sec:axioms}

We list the four paper-backed axioms that appear in
\texttt{\#print axioms Problem625.Publishable.erdos\_625\_full\_clean}, with
their literal sources and explicit disclosures. Two are literal
one-citation results from Heckel--Panagiotou~\cite{heckelpanagiotou2023};
two are hybrid (HP-2023 + auxiliary content).
See \texttt{proof/red-team/README.md} for the chronological index of all
audit artefacts.

(Note: the three ``layer-1'' axioms appearing in the legacy
\texttt{erdos\_625} theorem ---
\texttt{paperPartBThresholdAverage...source},
\texttt{paperPartBExactNoEmpty...source}, and
\texttt{heckel\_offdiag\_term\_bound} ---
are \emph{not} reachable from \texttt{erdos\_625\_full\_clean} because the
good case of the flagship goes through \texttt{erdos\_625\_97}, which uses
\texttt{lemma\_7\_20\_modified} instead. They are documented in the
repository's \texttt{paper/SOURCES.md}.)

\paragraph{Axiom numbering note.} We use a single axiom numbering
across the entire package (README, SOURCES, proof.md, this paper).
A1 = \texttt{lemma\_7\_20\_modified}; A2 =
\texttt{partB\_alphaMinusTwo\_firstMomentBelowOne\_source};
A3 = \texttt{chi\_alphaMinusTwo\_lower\_bound\_whp};
A4 = \texttt{zeta\_alphaMinusTwo\_upper\_bound\_whp}. A1 is
\emph{hybrid} (peer-reviewed bulk + in-repository numerical
certificate); A4 is an \emph{extrapolation} (from $(\alpha-1)$ to
$(\alpha-2)$); A2 and A3 are literal one-citation HP-2023 quotes.

\paragraph{Axiom~A1: \texttt{lemma\_7\_20\_modified}.}
(\emph{Hybrid:} HP-2023 Lemma~7.20 (modified) + in-repository
numerical certificate \texttt{lemma\_7\_10\_ext}.) Used on the
good-case branch (\emph{i.e.\@}\ under $\inrangemod(\varepsilon,n)$).
\begin{itemize}
\item The peer-reviewed bulk is HP-2023 Lemma~7.20, modified by
weakening condition (d) to $\mu_{\alpha}\geq n^{\xnaught+\varepsilon}$,
which Heckel~\cite{heckel2024} \S Discussion explicitly conjectures as
``straightforward''.
\item The numerical certificate filling the $\phi$-positivity gap
$[\xnaught+\varepsilon,0.04)$ outside HP-2023 Lemma~7.10's coverage is
ours (not peer-reviewed). The certificate is a $1086$-cell grid with a
Lipschitz envelope and a $10\%$ safety buffer on the Lipschitz
constant; positivity is established by envelope lower bound
$\geq 6.15\times 10^{-7}$ against a Lipschitz slack $3.74\times 10^{-8}$,
\emph{i.e.\@}\ $\approx 1.2$ orders of magnitude. Reproducibility: the
Python script
\texttt{proof/red-team/num-gap-lemma710-extension-2026-05-10.py}.
\end{itemize}
The natural follow-up project is either to publish the numerical lemma
as a standalone arXiv note, or to formalize the certificate in Lean
via \texttt{decide} on a rational grid. See
\texttt{proof/red-team/lemma-7-10-ext-disclosure-2026-05-11.md} for the
audit.

\paragraph{Axiom~A2: \texttt{partB\_alphaMinusTwo\_firstMomentBelowOne\_source}.}
(Pure paper citation.) HP-2023 Lemma~8.1 first-moment input at level
$\alpha-2$: for $k<\partbthr(n)+\lceil n/\log^{2}n\rceil$, the expected
number of $(\alpha-2)$-bounded proper $k$-colorings of $\Gnh$ is less
than $1$. The Lean theorem \texttt{kThreshold\_gap\_alpha\_minus\_2}
(proved with $0$ \texttt{sorry}) packages this into the deterministic
threshold gap of \Cref{lem:gap}.

\paragraph{Axiom~A3: \texttt{chi\_alphaMinusTwo\_lower\_bound\_whp}.}
(Pure paper citation, by analogy with HP-2023 \S 8.) \Cref{lem:chi_lower}:
$\chrom(G)\geq\athr(n)-n^{0.99}$ \whp\ $\geq 1-\varepsilon$, on the
crossing event $\neg\inrangemod(\varepsilon,n)$. The proof is the same
first-moment + Azuma argument that HP-2023~\S 8 uses for
$\chrom_{a}\geq\mathbf{k}_{a}-1$ \whp; the slack $n^{0.99}$ absorbs
both the residual additive $-1$ from Lemma~8.1's conclusion and the
Azuma--Hoeffding deviation. HP-2023~\S 8 itself applies this argument
at $a\in\{\alpha-1,\alpha-2\}$.

\paragraph{Axiom~A4: \texttt{zeta\_alphaMinusTwo\_upper\_bound\_whp}.}
(\emph{Extrapolation:} the symmetric $(\alpha-2)$-version of
Heckel~\cite{heckel2024} Proposition~5(b), obtained via HP-2023's
general-$a$ second-moment lemmas (Lemmas~6.3--6.5).)
\Cref{lem:zeta_upper}: $\cochrom(G)\leq\partbthr(n)+n^{0.99}$ \whp\
$\geq 1-\varepsilon$, on the crossing event. The proof transfers
Heckel~\cite{heckel2024} \S 5--6 (cochromatic Paley--Zygmund + Azuma)
from the published $(\alpha-1)$-bounded case to the $(\alpha-2)$-bounded
case. The transfer is principled but \emph{not literally in print}: the
underlying HP-2023 second-moment lemmas
(\cite{heckelpanagiotou2023}~Lemmas~6.3--6.5) are stated for general
$a$-bounded profiles, but the application step at $a=\alpha-2$ is not a
verbatim quote of any published lemma. We use the term ``extrapolation''
consistently across the package (README, SOURCES.md, proof.md) for this
axiom. The detailed audit is in
\texttt{proof/red-team/heckel2024-alpha-minus-two-transfer-audit-2026-05-11.md}.

\paragraph{Lean kernel axioms.}
The complete \texttt{\#print axioms} output of
\texttt{erdos\_625\_full\_clean} also lists \texttt{propext},
\texttt{Classical.choice}, \texttt{Quot.sound}. These are the standard
classical-logic axioms of Lean~4 / Mathlib4, present in essentially
every Mathlib proof.

\paragraph{Summary count.} Four paper-backed axioms (\emph{two}
literal: A2, A3; \emph{one} hybrid: A1; \emph{one} extrapolation: A4)
and three Lean kernel axioms --- seven entries total.

\section{Scope of the result}\label{sec:scope}

\paragraph{What we prove.}
\Cref{thm:main} delivers the in-probability divergence
$\chrom(\Gnh)-\cochrom(\Gnh)\to\infty$ for \emph{all} sufficiently
large $n$, with a quantitative rate $n^{1-2\varepsilon}$ for any
fixed $\varepsilon\in(0,0.001)$, with $0$ sorry on the Lean proof path
and seven axioms total (four paper-backed plus three Lean kernel).

\paragraph{What we do not yet prove (in Lean).}
\begin{enumerate}
\item \emph{Almost-sure convergence; \Cref{cor:in_prob}.} The
Erd\H{o}s--Gimbel question literally asks for ``almost surely''.
\Cref{thm:main} delivers a fixed-$\varepsilon$ in-probability bound;
the diagonal $\varepsilon_{n}=1/\log n$ Borel--Cantelli step that
promotes this to almost-sure convergence is a routine measure-theoretic
argument not yet in the Lean chain. Also, \Cref{cor:in_prob} stated in
\S 1 (in-probability divergence) is the
$\varepsilon\to 0$ limit of \Cref{thm:main}, which itself is a
straightforward exercise but is not in the Lean chain either; only
the per-$\varepsilon$ quantitative form is machine-checked.
\item \emph{Heckel's conjectured sharp rate.} Heckel conjectures
$\chrom-\cochrom=\Theta(n/\log^{3}n)$ \whp; our $n^{1-2\varepsilon}$ is
strictly weaker.
\item \emph{Axiom-free Lean proof.} All four paper-backed axioms are
defended as the natural unit of citation, with explicit disclosure for
the hybrid (axiom~4) and the symmetric extension
(\Cref{lem:zeta_upper}). A fully discharged proof is a 12--24~months
formalization project (the HP-2023 tameness + second-moment chain,
plus the numerical certificate of axiom~4).
\end{enumerate}

\section*{Acknowledgments and provenance}

\paragraph{Provenance disclosure.} The Lean formalization, this
paper, the proof writeup \texttt{proof/proof.md}, and the five
internal red-team audit passes in \texttt{proof/red-team/} were
generated by an LLM-agent pipeline (Anthropic Claude Opus~4.7,
model id \texttt{claude-opus-4-7[1m]}, the 1M-context variant, driven
by the \texttt{operator} agent orchestrator at
\url{https://github.com/uthunderbird/vibechord}, whose codex-brain
adapter ran \texttt{gpt-5.3-codex-spark} at \texttt{effort=low})
under human-in-the-loop supervision. The 1M-context variant was
load-bearing: \texttt{PartBProfileBridge.lean} is 29512 lines
($\sim$1.59~MB). The
human supervisor selected the proof strategy, accepted or rejected
proposed proof steps, and ran the Lean build; the LLM agents
generated and edited the Lean source, the writeups, and the
adversarial-audit artefacts. The five \texttt{/swarm-red-team}
sessions referenced below were conducted by the same LLM-agent
pipeline that produced the proof and are therefore \emph{internal
adversarial audits}, not independent third-party reviews. We disclose
this prominently because it is a genuine novelty of the work and
because a reader scoring the contribution should know which
correctness signal is machine-checked (the Lean kernel verifies the
proof modulo the seven axioms enumerated in \S\ref{sec:axioms}) and
which is not (the writeups, the axiom-to-paper correspondence, the
hybrid-axiom disclosures, and the audit reports rely on the same
LLM-agent pipeline that produced them).

\paragraph{Acknowledgments.} We thank Annika Heckel and Konstantinos
Panagiotou for the foundational papers
\cite{heckel2024,heckelpanagiotou2023} that this formalization builds
on, and the Lean community for the Mathlib library. The complete
Lean development, red-team artefacts, and numerical scripts are
publicly available at the repository link listed in the abstract.
The development was adversarially audited internally via five
\texttt{/swarm-red-team} sessions (2026-05-11 on the Lean theorem,
2026-05-12 on \texttt{proof.md}, on the present paper, and on a
final strict lemma-by-lemma pass); each session reported no critical
(P0) issues in its respective round. As above, these are internal
audits of an LLM-agent pipeline by the same pipeline; readers should
treat them as an internal-quality-gate signal, not as third-party
review. Agent-side reproducibility inputs (operator profile,
harness instructions, model/effort settings) are summarized in
\texttt{README.md} and \texttt{DEVELOPMENT.md}; machine-local harness
dumps are intentionally not shipped in the publication package.

\paragraph{Limitation framing of LLM-generated mathematics.} Readers
should be aware that LLM-generated proofs may exhibit
plausibility-driven failure modes that can survive same-model
internal critique because the proof author and the critic share
training-data blind spots. Here the proof author and the five
internal red-team critics are the same model
(\texttt{claude-opus-4-7[1m]}); the internal audits are therefore
structurally incapable of detecting blind spots inherited from
training. External verification by an unrelated reader, a different
model, or a human mathematician is encouraged. The Lean kernel
acceptance (build GREEN, \texttt{\#print axioms} machine-checked)
reduces but does not eliminate this risk: the four paper-backed
axioms remain admitted, and the natural-language disclosures
surrounding them --- axiom-to-paper correspondence, the hybrid
disclosure for A1, and the extrapolation disclosure for A4 ---
remain LLM-authored. Cross-model auditing was technically available
through the \texttt{operator} framework's adapter layer but was not
used in this development; this is the obvious mitigation the
pipeline did not take.

\paragraph{Anthropic non-endorsement.} Anthropic is acknowledged as
the provider of the Claude model used. No endorsement of this work
by Anthropic, OpenAI, or any other organization is implied.

\paragraph{Human supervisor.} Daniyar Supiyev directed the proof
strategy at the branch-point level, accepted/rejected proposed proof
steps, ran the Lean build, and chose which red-team findings to act
on. Contact is via GitHub issues at the repository linked in the
abstract; the supervisor's role is non-anonymous for this release.
No institutional affiliation is claimed.

\begin{thebibliography}{9}

\bibitem{erdosproblems625}
The Erd\H{o}s Problems. Problem 625. \\
\url{https://www.erdosproblems.com/625}.

\bibitem{heckel2024}
A.~Heckel,
\textit{The difference between the chromatic and the cochromatic number
of a random graph}. arXiv:2409.17614, 2024.

\bibitem{heckelpanagiotou2023}
A.~Heckel and K.~Panagiotou,
\textit{Colouring random graphs: Tame colourings}. arXiv:2306.07253, 2023.

\bibitem{heckelriordan2023}
A.~Heckel and O.~Riordan,
\textit{How does the chromatic number of a random graph vary?}
arXiv:2103.14014, 2021 (v3 updated 2023).

\end{thebibliography}

\end{document}
